\documentclass{resume}
\usepackage{zh_CN-Adobefonts_external} 
\usepackage{linespacing_fix}
\usepackage{cite}
\usepackage{enumitem}
\usepackage{url}
\usepackage{hyperref}
\hypersetup{colorlinks=true,linkcolor=black,urlcolor=black,citecolor=black}
% ============ 页眉改版：左对齐 + 照片固定右上角 ============
\RequirePackage{eso-pic}
\makeatletter

% 姓名：左上角大字
\renewcommand{\name}[1]{%
  {\Huge\scshape #1}\par\vspace{0.4ex}%
}

% 联系方式：手机 · 邮箱 同一行，左对齐
\renewcommand{\contactInfo}[2]{%
  {\sffamily\large 手机：#1 \textperiodcentered\ 邮箱：#2}\par\vspace{0.4ex}%
}


% 附加信息：左对齐，空参数自动省略分隔符
\renewcommand{\otherInfo}[4]{%
  {\sffamily\large #1
   \ifthenelse{\isempty{#2}}{}{\ \textperiodcentered\ #2}%
   \ifthenelse{\isempty{#3}}{}{\ \textperiodcentered\ #3}%
   \ifthenelse{\isempty{#4}}{}{\ \textperiodcentered\ #4}}\par\vspace{0.3ex}%
}

% 照片：固定右上角（右缘贴 0.7in 右边距，顶端贴 0.5in 上边距，与 geometry 一致）
\renewcommand{\yourphoto}[1]{%
  \AddToShipoutPictureBG*{%
    \AtPageUpperLeft{%
      \hspace{\dimexpr\paperwidth-0.7in-#1\paperwidth\relax}%
      \setbox\@tempboxa=\hbox{\includegraphics[width=#1\paperwidth]{images/you.jpg}}%
      \raisebox{-\dimexpr0.5in+\ht\@tempboxa\relax}{\box\@tempboxa}%
    }%
  }%
}
\makeatother

% 论文条目：标题整段自然换行（不连字符断词、不超行宽）；日期写在署名/状态行右侧
\newcommand{\papertitle}[1]{%
  \par\noindent{\normalsize\hyphenpenalty=10000\exhyphenpenalty=10000\textbf{#1}}\par
}

\begin{document}
\pagenumbering{gobble}

% ************************** 基础信息 **************************
\name{李子林}
\contactInfo{(+86) 19980528634}{2023141461067@stu.scu.edu.cn}
\otherInfo{\textbf{四川大学} · \textbf{软件工程专业} · \textit{本科2023级}}{}{}{}
\yourphoto{0.06}


% ************************** 学术成果（保研核心！置顶突出） **************************
\section{\textbf{科研经历}}
{\small
\papertitle{\textbf{BRACE: Taming Sharp Irregularities via Barycentric Rational Forecasting for Fast Diffusion Transformers Inference}}
\begin{itemize}[label=\textbullet, parsep=0.2ex, itemsep=0.1ex]
  \item 面向扩散 Transformer 推理开销大、现有长步长预测稳定性不足的问题，提出无需训练的 BRACE 加速框架；通过局部滑动窗口、重心有理外推与自适应 Chebyshev 权重实现中间特征稳定预测，在 DiT-XL/2、FLUX.1-dev 与 HunyuanVideo 上均取得优于现有方法的质量–效率权衡。
  \item 作为第二作者被 ACM MM 2026 接收 (Poster)。\hfill \mbox{2025.11 – 2026.4}
\end{itemize}

\papertitle{\textbf{HDS: Heterogeneity-Driven Safeguarding for Accelerating Diffusion Transformers}}
\begin{itemize}[label=\textbullet, parsep=0.2ex, itemsep=0.1ex]
  \item 面向现有预测加速方法忽略误差影响跨层异质性、难以兼顾质量与效率的问题，提出无需训练的 HDS 加速框架：基于预测误差影响具有时间稳定层间异质性的发现，设计异质性分数（HS）在线评估各层风险，对脆弱层保留完整计算、对稳健层激进加速，在多种 DiT 架构与预测方案上以可忽略开销提升质量–效率帕累托前沿。
  \item 作为第一作者，投稿 CVPR 2027 (Under Review)。\hfill \mbox{2026.03 – 2026.05}
\end{itemize}

\papertitle{\textbf{ExProp: Exponential Propagation on Implicit Manifolds for Training-Free DiT Acceleration}}
\begin{itemize}[label=\textbullet, parsep=0.2ex, itemsep=0.1ex]
  \item 面向现有特征预测方法仅按欧氏空间轨迹外推、忽视扩散轨迹隐流形结构的问题，提出无需训练的 ExProp 加速方法：将特征预测重构为隐流形上的局部指数传播，通过局部信任域机制约束几何有效邻域，在 Flux.1-dev、DiT-XL/2 与 HunyuanVideo 等主流模型上优于现有特征缓存与预测方法。
  \item 作为第一作者，投稿 CVPR 2027 (Under Review)。\hfill \mbox{2026.05 – 2026.07}
\end{itemize}

\papertitle{\textbf{FADE: Fractional Anomalous Dynamics Extrapolation for Training-Free Diffusion Transformers Acceleration}}
\begin{itemize}[label=\textbullet, parsep=0.2ex, itemsep=0.1ex]
  \item 面向扩散 Transformer 潜在轨迹非平稳、高曲率导致整数阶预测算子误差累积的问题，提出无需训练的 FADE 加速框架：以反常扩散刻画特征演化，采用 Mittag–Leffler 分数阶算子适配高曲率非线性轨迹，并通过离线多维建模实现推理期近零开销的自适应参数检索，在多种架构与模态上实现领先的加速与保真度权衡。
  \item 作为第二作者，被 NeurIPS 2026 接收。\hfill \mbox{2026.03 – 2026.05}
\end{itemize}

\papertitle{\textbf{LaBACK: From Feature Forecasting to Trajectory-Drift Control for Efficient Flow Matching Sampling}}
\begin{itemize}[label=\textbullet, parsep=0.2ex, itemsep=0.1ex]
  \item 面向流匹配采样中缓存引入的速度误差沿轨迹传播造成生成漂移的问题，提出无需训练的 LaBACK 加速方法：以 Laurent 级数引导的特征预测降低近似误差，并引入确定性回溯抑制误差沿采样轨迹传播，在文生图、蒸馏文生图与文生视频任务上取得更优的速度–质量权衡。
  \item 作为第二作者，投稿 CVPR 2027 (Under Review)。\hfill \mbox{2026.05 – 2026.07}
\end{itemize}
}

\vspace{0.05em}

% ************************** 项目经历 **************************
\section{\textbf{项目经历}}
\datedsubsection{\textbf{SEDD-LaTeX-RLHF | 离散扩散 RLHF 对齐}}{2025.09}
\begin{itemize}[parsep=0.2ex, itemsep=0.1ex]
  \item 将 PPO 算法适配到离散扩散架构（SEDD），完成 SFT + PPO 强化学习对齐训练。
  \item 通过稀疏采样优化将显存从 GiB 级降至 MB 级。
  \item 设计多维 reward 函数评估 LaTeX 语法与数学质量，RL 对齐后生成质量显著优于 SFT 基线
  
  \url{https://github.com/zlven/SEDD-LaTeX-RLHF}
\end{itemize}

\vspace{0.05em}

\datedsubsection{\textbf{Medical Grand Challenge（NUS） | CT-MRI 跨模态 3D 建模系统}}{2025.10 - 2026.08}
\begin{itemize}[parsep=0.2ex, itemsep=0.1ex]
  \item 基于 3D Slicer 构建 CT-MRI 跨模态3D建模系统，用于精确外科手术的术前诊断规划。
  \item 通过多阶段互增强架构（BSAFusion + 跨模态注意力），优化跨模态对齐与小病灶分割精度。
  \item 采用KV cache与 int8 量化轻量化，将耗时单位由天降至分钟级别、显存占用减少50\%，以适配临床常规工作站批量处理。
\end{itemize}
% ************************** 实习经历 **************************
\section{\textbf{实习经历}}
\datedsubsection{\textbf{山西省智慧交通实验室有限公司（信息技术研究院）}  \textbf{算法工程师}}{2026.07 - 2026.08}
\begin{itemize}[parsep=0.2ex, itemsep=0.1ex]
  \item 参与黑光相机低光照场景下计算机视觉检测模型的训练与调优，面向弱光/夜间小目标检测，优化多尺度特征表达与模型结构，结合模型压缩实现轻量化部署，提升目标检测与识别性能。
  \item 参与撰写《改扩建工地无人机视觉违规识别与远程喊话干预系统》发明专利（在审），包括算法设计、可行性分析等。
  
\end{itemize}

% ************************** 个人荣誉 **************************
\section{\textbf{个人荣誉}}
\begin{itemize}[parsep=0.2ex, itemsep=0.1ex]
  \item 2024全国大学生数学建模竞赛 \textbf{四川省一等奖}
  \item 2025全国大学生数学建模竞赛 \textbf{四川省二等奖}
  \item 2025中国高校计算机大赛-AIGC创新赛 \textbf{西南赛区三等奖}(软著1)
  \item 2026 Medical Grand Challenge（NUS）\textbf{Open 赛道第四名}
  \item 四川大学优秀学生干部
\end{itemize}

\vspace{0.05em}


% ************************** 专业技能 **************************
\section{\textbf{专业技能}}
\begin{itemize}[parsep=0.2ex, itemsep=0.1ex]
  \item \textbf{技术栈}: Python、PyTorch、Java、C/C++、Linux、 git等
  \item 掌握操作系统、计算机网络、数据结构与算法，深度学习等相关基础知识
  \item \textbf{英语水平}: CET-4 、CET-6
\end{itemize}

\end{document}